% Generated draft: requires mathematical and source review.
\documentclass[11pt]{article}
\usepackage[margin=25mm]{geometry}
\usepackage{fontspec}
\setmainfont{Latin Modern Roman}
\usepackage{amsmath,amssymb,mathtools,unicode-math}
\setmathfont{Latin Modern Math}
\usepackage{enumitem,needspace,xurl,hyperref}
\setlength{\parindent}{0pt}
\setlength{\parskip}{5pt}
\setlength{\emergencystretch}{3em}
\newcommand{\N}{\mathbb{N}}
\newcommand{\Z}{\mathbb{Z}}
\newcommand{\Q}{\mathbb{Q}}
\newcommand{\R}{\mathbb{R}}
\newcommand{\C}{\mathbb{C}}
\newcommand{\F}{\mathbb{F}}
\newcommand{\A}{\mathbb{A}}
\newcommand{\PP}{\mathbb{P}}
\newcommand{\E}{\mathbb{E}}
\newcommand{\eps}{\varepsilon}
\newcommand{\dd}{\,\mathrm{d}}
\begin{document}
\section*{UnsolvedMath ID: 30004502}
\textbf{Source problem code:} catalogue entry
\subsection{Definitions and mathematical statement}
Let \(\zeta(s)\) denote the Riemann zeta function, let \(i^2=-1\), and let \(\log\) denote the natural logarithm. For \(0\leq q\leq1\), define
\[
M_q(T)=\frac1T\int_T^{2T}\left(\int_{-1/2}^{1/2}\left|\zeta\!\left(\tfrac12+it+ih\right)\right|^2\,\dd h\right)^q\,\dd t.
\]
The first supplied estimate is
\[
M_q(T)\ll A_q(T),\qquad
A_q(T)=\left(\frac{\log T}{1+(1-q)\sqrt{\log\log T}}\right)^q.
\]
The natural asymptotic interpretation is \(T\to\infty\); the record gives no explicit lower threshold for \(T\). The displayed expression is real-valued for \(T\geq e\).

For a prime \(r\), let \(\chi\) range over Dirichlet characters modulo \(r\), and let \(\chi_0\) be the principal character. For the supplied range \(1\leq x\leq r\), put
\[
L=\min\{x,r/x\},\qquad S_\chi(x)=\sum_{\substack{n\in\Z\\1\leq n\leq x}}\chi(n).
\]
For \(r>2\) and \(0<q\leq1\), define
\[
C_q(r,x)=\frac1{r-2}\sum_{\substack{\chi\ ({\rm mod}\ r)\\\chi\ne\chi_0}}|S_\chi(x)|^{2q}.
\]
The second supplied estimate is
\[
C_q(r,x)\ll B_q(r,x),\qquad
B_q(r,x)=\left(\frac{x}{1+(1-q)\sqrt{\log\log(10L)}}\right)^q.
\]
Here \(U\ll V\) means that \(U\leq K V\) for a positive constant \(K\); which parameters \(K\) may depend on is not specified in the record. The source does not specify whether \(x\) is real or integer; the displayed sum has a meaning in either case.

\textbf{Supplied sharpness conjecture, attributed to Harper.} Are both displayed upper estimates sharp for every \(0\leq q\leq1\), including their displayed logarithmic factors? An order-of-magnitude interpretation would ask for corresponding lower estimates \(M_q(T)\gg A_q(T)\) and \(C_q(r,x)\gg B_q(r,x)\) in the intended parameter regimes. The record does not specify whether these lower estimates should hold throughout the ranges, along suitable sequences, or with constants uniform in \(q\).

\textbf{Unresolved qualifications.} A lower estimate throughout the stated character range cannot be the literal intended assertion: at \(x=r\), every nonprincipal character has \(S_\chi(r)=0\), whereas \(B_q(r,r)>0\) for \(q>0\). No replacement range or normalization is supplied. Also, the stated allowance of every prime includes \(r=2\), where the normalization \(1/(r-2)\) is undefined; the definition above is consequently restricted to the meaningful case \(r>2\), pending clarification. At \(q=0\), vanishing character sums produce \(0^0\), for which no convention is supplied. Thus the full endpoint-inclusive conjecture requires clarification before it becomes a precise two-sided assertion.
\subsection{Short English statement}
The supplied conjecture, attributed to Harper, asks whether two moment upper bounds have the correct order, including their logarithmic factors, for every \(0\leq q\leq1\). The intended meaning of sharpness and uniformity is unspecified. The character formulation also needs clarification at \(r=2\), at \(q=0\), and near the endpoint \(x=r\), where a matching positive lower bound fails.
\subsection{Sources}
\begin{itemize}
\item Local UnsolvedMath registry: permanent ID \(30004502\), title ``Sharp Moments of Zeta and Character Averages,'' exact source problem code \texttt{catalogue entry}, in the Oberwolfach Reports Open Problems collection. The original statement attributes the conjecture to Harper; the separate proposer field is empty. The supplied proposed year is 2020.
\item Supplied report citation: \emph{Automorphic Forms and Arithmetic} (2021), Oberwolfach Reports, pp.\ 1317--1383, DOI \texttt{10.4171/owr/2020/26}. Supplied source URL and literature reference ``Oberwolfach Report 26/2020'': \url{https://doi.org/10.4171/owr/2020/26}. The citation's 2021 date and the report's 2020 designation are retained as supplied.
\item Supplied literature assessment, dated exactly \texttt{2026-08-22}: the record labels the problem open, states that the source explicitly presents the sharpness assertion as a conjecture, and reports that no primary result establishing both sharp bounds was found. This is supplied dated curation, not an independently verified literature assessment. The background explicitly requests independent verification of time-sensitive status and cited claims.
\item Supplied statement-verification metadata labels the statement \texttt{corrected\_verified} and reports that the official report was checked and both moment bounds and their parameters were restored. This describes the supplied provenance; no source visit or independent verification is claimed here.
\end{itemize}
\end{document}
